\documentclass[11pt]{article}
\usepackage[margin=1in]{geometry}
\usepackage{amsmath,amssymb,hyperref,xcolor}
\newcommand{\K}{\mathrm K}
\newcommand{\RES}{\mathrm{RES}}
\newcommand{\RV}{\mathrm{RV}}
\newcommand{\FN}{\mathrm{FN}}
\newcommand{\Gm}{\mathbb G_m}
\title{Section 4: first mathematical review}
\author{Proposed passages for author review}
\date{November 11, 2026}
\begin{document}
\maketitle
\noindent\textbf{Status.} Historical supporting document: the review destination
is now the working manuscript, where the proposed passages appear in dark
orange awaiting author review. Source: repository checkpoint \texttt{82fe1a1}, original annotated
PDF pp.~36--43. Main-text statements were read independently before the
annotations. The accompanying Lean file checks the algebraic and
ordered-set calculations below; it does not yet prove the numbered
retractions or their concrete Grothendieck constructions.

\section*{1. Cell dimensions and homogeneous notation (R056--R059)}
\textbf{Remark 4.2, proposed replacement for the cube description.}
Choose a finite cell decomposition $X=\coprod_j C_j$ and set
\[
  \chi_g(X)=\sum_j(-1)^{\dim C_j}.
\]
Equivalently, a rectilinearization uses open boxes of varying dimensions,
including dimension zero. A singleton already requires the zero-dimensional
case. Independence of decomposition, invariance and multiplicativity remain
proof obligations; this correction is not a substitute for those proofs.

\medskip\noindent\textbf{Notation 4.3, proposed clarification.}
In the graded ring write $u=[\{1\}]_1$ and
$L_1=[\Gm]_1+u$; the ring unit is $1=[\{1\}]_0$.
Graded localizations by $[\mathbb A]$ mean localization by $L_1$.
The ideal $!I[*]$ is generated by the homogeneous degree-one differences
$[\gamma^\sharp]_1-[\Gm]_1$. In the ungraded ring the corresponding
affine-line class is $L=[\Gm]+1$.

\medskip\noindent\textbf{Display (4.2), proposed replacement.}
With $T=[\Gm]_1$, the calculation must read
\[
 \mathbb E_g([\RV^{\circ\circ}]+1-u)
   =-T L_1^{-1}+1-uL_1^{-1}=0.
\]
Indeed $L_1=T+u$. The printed last term $-L_1^{-1}$ omits the
degree-one point factor; its algebraic remainder is $(u-1)L_1^{-1}$.
After forgetting degree, the point is $1$ and the familiar ungraded
expression is recovered. The bounded version uses
$1-u u^{-1}=0$. Component tensor maps are first bilinear over
$\mathbb Z$; their sum must then respect the full homogeneous balancing
relations (R058). This last step is a formalization obligation, not a
change of target.

\section*{2. Degree-zero constants (R062)}
\textbf{Remark 4.6, proposed replacement.}
Let
\[
 D=\Gamma\cap\operatorname{dcl}^{\mathrm{eq}}(\mathbb S),
 \qquad \mathbb Z[D]\subseteq\FN(\Gamma,\mathbb Z)
\]
be the group ring spanned by the $p_\gamma$ for all definable Gamma
elements. Then $\K\mu\Gamma[0]\cong\mathbb Z[D]$, and the target
of $\lambda$ is $\mathbb Z[D]\oplus X\FN(\Gamma,\mathbb Z)[X]$.

\noindent\textbf{Reason.} A degree-zero object is a finite collection of
empty tuples with definable constant weights. A singleton of weight
$\tfrac12 v(t)$ over $\mathbb C((t))$ is permitted even though this value
is not in the parameter field's value group $\mathbb Z v(t)$.
The earlier remark after the Gamma-category definition already distinguishes
these two groups. This proposal makes the coefficient notation consistent
with that remark. All three assignments for $\psi$ keep their stated
degree restrictions: only $p_\gamma$ occurs in degree zero (R063).

\section*{3. The sign in Proposition 4.7 (R064)}
\textbf{Proposed third bullet.}
\[
 \mu\mathbb E_g(\mathfrak h)=+[1]/([\mathbb A]),
 \qquad \mu\mathbb E_b(\mathfrak h)=0.
\]
\textbf{Derivation from the paper's definitions.} For
$\boldsymbol I=((0,\infty),-\mathrm{id})$, the corrected weight
$\mu_I(x)=-x+x$ is zero. Thus
\[
 \lambda_g([\boldsymbol I])=-p_0X,
 \qquad \psi(-p_0X)=-[\Gm]_1=+[1]_1
       \pmod{([\mathbb A])}.
\]
The last equality is forced by $[\mathbb A]_1=[\Gm]_1+[1]_1$.
For the bounded characteristic the open ray has characteristic zero.
The original negative sign is not the value of this construction.
This is a change to a proposition's stated value and needs author approval.
It also requires checking later callers, including the use of 4.7 in \S8.

\section*{4. Bounded support and negative intervals (R067--R068)}
\textbf{Notation 4.11, proposed replacement.}
Define $\FN^{\mathrm{db}}(\Gamma,\mathbb Z)$ to consist of the functions
in $\FN(\Gamma,\mathbb Z)$ with doubly bounded nonzero locus. Set
\[
 o_\gamma=\begin{cases}
 \mathbf1_{(0,\gamma)},&\gamma>0,\\
 \mathbf1_{(\gamma,0)},&\gamma<0,\\
 -p_0,&\gamma=0.
 \end{cases}
\]
The existing two alternative convolution formulas are then read with
explicit sign cases. In particular
\[
 o_\alpha*o_\beta=
 \begin{cases}
 -o_{\alpha+\beta},&\alpha,\beta\geq0\text{ or }\alpha,\beta\leq0,\\
 -o_\alpha-o_\beta-p_0,&\alpha<0<\beta\text{ or }\beta<0<\alpha.
 \end{cases}
\]
For $\beta\ne0$, $p_\alpha*o_\beta$ is the characteristic function
of the open interval with endpoints $\alpha$ and $\alpha+\beta$;
for $\beta=0$ it equals $-p_\alpha$.

\noindent\textbf{Reason.} Under the literal printed convention,
$\gamma<0$ gives the empty interval $(0,\gamma)$ and hence $o_\gamma=0$.
The proposed relation $o_\gamma+p_0=0$ would then kill $p_0$, the
convolution unit, contradicting the claimed quotient $\mathbb Z$.
The special convention $o_0=-p_0$ concerns functions, not the actual
empty interval $(0,0)$.

\section*{5. Admissible radii and the zero generator (R069--R071)}
\textbf{Notation 4.13, proposed replacement.}
Define $\boldsymbol P_\gamma$ for nonnegative definable $\gamma$
\emph{admitting a definable} $t_\gamma\in\gamma^\sharp$. Explicitly,
the index set is
\[
 \{\gamma\geq0:\ \exists t\in\RV(\operatorname{dcl}\mathbb S),\
                  \operatorname{vrv}(t)=\gamma\}.
\]
Use this same index set for the generators of $(\boldsymbol P_\Gamma)$.
This preserves the intended definition wherever its specified point exists.
It does not assume a definable section of the value map.

\noindent\textbf{Reason.} The distinction is already explicit after the
definition of Gamma-categories. Over $\mathbb C((t))$, the value
$v(t)/2$ is definable, but its RV fibre has no definable point.
For example, the conjugation sending $t^{1/2}$ to $-t^{1/2}$ changes
its RV class. Such a radius cannot be used to select $t_\gamma$ over
the unchanged base.

\medskip\noindent\textbf{Proposition 4.15, proposed proof replacement at
the generator calculation.} If $\gamma=0$, the annulus is empty and
$[\{t_0\}]-[1]=0$ by the zero-volume, value-zero singleton isomorphism;
thus $\boldsymbol P_0=0$. If $\gamma>0$, its Gamma part is
$(0,\gamma]\,=(0,\gamma)\amalg\{\gamma\}$ and
\[
 \chi((0,\gamma))[\Gm]_1+
 \chi(\{\gamma\})[\Gm]_1+[1]_1-[1]_1
 =-[\Gm]_1+[\Gm]_1=0.
\]
The displayed open-interval-plus-endpoint expansion in the original proof
does not apply at zero; there it leaves an extra $[\Gm]_1$.
This is a proof correction; the intended bounded retraction is retained.

\section*{6. Matters not approved or completed by this review}
The affine-localized injectivity assertion in Remark 4.4 (R060) has a
separate independent paper-level reconstruction in the companion file
\texttt{R060-independent-review.md}, checked against the primary source.
The generic localization mechanism is checked in Lean; the actual geometric
and graded identifications are not yet formalized.
The assertion has not been silently replaced by ``natural homomorphism.''
The point-localized comparison, convolution associativity (R061), enlarged
alternative quotient (R066), and diagram shorthand (R072) still need their
actual formal constructions; they are not assumed false on account of
annotations. Pure editorial changes R055, R065, R070 and R168 can be
collected separately. No referee annotation has been marked resolved yet.

\medskip\noindent\textbf{Author decision requested.} Approve the proposed
cell/homogeneity clarifications and mathematical corrections in items
1--5, or specify which passage should retain a different interpretation.
The highest-impact decision is the sign in Proposition 4.7. Full coverage,
remaining proof obligations and exact formalization scope are recorded in
\texttt{section 4/COVERAGE.md}.
\end{document}
